\section{Empirical Objects}

The empirical objects are built to match the theory as closely as the local data allow.
The raw investable universe is filtered to liquid, finite-Greek listed options on the
eight primary single-name underlyings. Because listed options expire, the empirical assets
are rolling option buckets rather than permanent CUSIPs or contracts. This is a necessary
measurement convention, not a change in theory. At each monthly snapshot the code selects
the liquid representative contract in each underlying/right/moneyness bucket, records its
mark and Greeks, and then computes the next holding-period bucket return from that
same prior-date choice. The local panel contains month-end snapshots of 15--21 DTE
options, so the selected option usually expires before the next snapshot. The realized
bucket return is therefore a listed-expiry payoff return: the payoff date is the listed
expiry date, and the terminal spot is the raw daily underlying close on that expiry. If a
split or unit change occurs between the decision date and expiry, the terminal raw close
is converted back into the decision-date contract units before it is compared with the
decision-date strike. The return is then
\[
  R_{i,t+1}=\frac{\max(S_{\tau}-K_i,0)}{C_{i,t}}-1
  \quad\text{or}\quad
  R_{i,t+1}=\frac{\max(K_i-S_{\tau},0)}{C_{i,t}}-1 .
\]
When a strategy assigns weight \(q_i\), the contribution to portfolio return is
\(q_iR_{i,t+1}\). The paid or received option premium is therefore embedded in both the
denominator and the NAV weight. The future expiry close is not a trading-decision
look-ahead: the option is selected at the decision date, and the future spot is used only
to realize the subsequent return after the listed holding period.

Four bucket types are used: ATM, near put, near call, and put wing. The construction
therefore contains both calls and puts by design. It does not start from a long-call
strategy. The optimizer may choose long or short signed exposures subject to gross NAV,
short-premium, per-contract, net NAV, and underlying concentration budgets. The empirical
weights are therefore option premium exposures, not share counts and not delta-equivalent
stock weights. The short-premium budget is not a transaction cost; it is a capital
accounting restriction that prevents a premium-only gross budget from becoming an
unbounded short-option tail exposure.

The risk model has three empirical components. First, underlying stock returns estimate
the covariance of spot shocks. Second, ATM implied-volatility changes estimate the
covariance of volatility shocks. Third, the residual option-return covariance is estimated
from the difference between realized bucket returns and the fitted Greek-factor return in
the training window. The residual is not a nuisance term to be optimized away. It is the
part of option risk that the local Greek map does not explain.

The expected-return vector is the training-window mean of option-bucket returns. This is
the weakest statistical input in the entire system, as it is in ordinary Markowitz. The
paper keeps it simple because the current objective is to validate the option-only
covariance and optimization mechanics. Forecasting expected option returns from volatility
risk premia, skew, dealer inventory, earnings events, and macro states is a separate
research problem.

The benchmark set is designed to make the option result harder to overstate. Equal premium
and equal risk are simple option-only baselines. Delta neutral is a constrained
option-only baseline. Random feasible portfolios test whether the optimizer is
doing better than arbitrary gross-NAV feasible option books. A delta-matched equity
portfolio holds the same first-order underlying delta vector as the Greek Markowitz option
book. An underlying-only Markowitz portfolio applies the same gross-NAV tangency logic
directly to the eight equities. These two equity benchmarks ask whether the option result
is merely an equity result in disguise.

All empirical returns are monthly. The train window ends on 2020-12-31. The test window
uses observations strictly after that date; the first test return is selected on
2020-12-31 and realized at the next snapshot. Bootstrap Sharpe intervals in the
performance table are computed from the realized test-window return series and are used as
uncertainty diagnostics, not as a license to search across many parameterizations.
Sortino and Omega use a zero monthly threshold, Calmar uses maximum drawdown, and the
information ratio uses SPY as the benchmark.
